\SourceLeaf{82}
\noindent\hspace*{\parindent}Es sei $\mathfrak{P}$ die verlangte Function, so ist:

\[
\begin{aligned}\left(\frac{d\mathfrak{P}}{dP}\right) ={}& \text{pars absoluta sive a } q \text{ libera in }\\ &\frac{q\,.\,f'q\,.\,f\dfrac{P^{\frac1n}}{Q}}{nP\left(1-\left(\dfrac{P}{Q}\right)^{\frac{\nu}{n}}\right)}.
\end{aligned}
\]



\Needspace{28\baselineskip}\begin{center}
BEMERKUNGEN ZU DEN NOTIZEN \"UBER REIHENINVERSION.
\end{center}

\noindent\hspace*{\parindent}Die erste Notiz ist aus einem Briefe an \textsc{Olbers} vom 12.\,Dec. 1813 entnommen. Der daselbst symbolisch durch $(n,\nu)$ bezeichnete Entwicklungscoefficient stellt sich auf Grund des polynomischen Lehrsatzes explicit dar in der Gestalt:

\[
(n,\nu) = \sum_\lambda \frac{(n+\mu-1)(n+\mu-2).\,.\,.(n+1)n}{\lambda'!\,\lambda''!\,\lambda'''!\,.\,.\,.\,.}\, a^{\prime\lambda'}a^{\prime\prime\lambda''}a^{\prime\prime\prime\lambda'''}.\,.\,.,
\]

\noindent wenn hierbei $a'=-\alpha$, $a''=-\beta$, $a'''=-\gamma$, .\,.\,. gesetzt wird; die Summe bezieht sich auf alle Systeme ganzer, nicht-negativer Zahlen $\lambda'$, $\lambda''$, .\,.\,., welche die Gleichung $\lambda'+2\lambda''+3\lambda'''+.\,.\,.=\nu$ befriedigen, und es ist $\mu=\lambda'+\lambda''+\lambda'''+.\,.\,.$.

Die zweite Notiz, welche \textsc{Gauss} in sein Exemplar von »\textsc{Hobert} und \textsc{Ideler}, trigonometrische Tafeln« (erste Seite des Umschlages) geschrieben hat, giebt in expliciter Gestalt die am Schlusse der ersten Notiz angedeutete Verallgemeinerung. Man wird die fragliche Inversionsformel mit H\"ulfe des eben entwickelten Ausdrucks f\"ur $(n,\nu)$ leicht best\"atigen.

In der dritten Notiz, welche sich auf einem einzelnen Zettel findet, und welche das Datum 1820 Jun.\,22 tr\"agt, ist nach einer geringf\"ugig abge\"anderten Methode gleichfalls die Verallgemeinerung der ersten Notiz f\"ur beliebige Potenzen von $x$ explicit entwickelt. Dem Schlusse von den beiden Darstellungen des $\Omega$ auf die Gleichungen I, II, .\,.\,. liegt eine Homogeneit\"atsbetrachtung zu Grunde.

Die vierte Entwicklung stammt vermuthlich aus dem Jahre 1816 und findet sich in einem Handbuche. Unter dem Symbol $[Fx]_{\text{Coeff.}\,x^n}$ ist hier der Coefficient von $x^n$ in der Entwicklung von $F(x)$ nach Potenzen von $x$ verstanden. \hfill\textsc{Fricke.}

